\section{Numerical illustration of the oracle comparison}
\label{sec:predictive-intro-figure}

Figure~\ref{fig:conference-tuned-intro} illustrates the two signal families in
Theorem~\ref{thm:conference-predictive} at $q=10^{-4}$ and $\sigma^2=1$. We
set $(\beta_C^*,\beta_R^*)=(0.1,1)$ for the rare-signal panel and $(1,0.1)$ for
the common-signal panel. All other assumptions match the theorem. We evaluate
61 integer observation caps, evenly spaced up to rounding from $10{,}002$ to
$40{,}004$, so $1\leq Mq(1-q)\leq4$. The horizontal axis reports the expected
number of rare observations at the cap, $Mq$, rather than a realized count.

For each rule, we use the parameterization in
\eqref{eq:conference-oracle-rates} and a grid consisting of zero and 4097
geometrically spaced positive values of $\alpha$ from $10^{-6}$ to $2$. We
iterate the exact recurrence \eqref{eq:predictive-moments}, including both
sources of training variance. At each cap, we select the rate and integer
stopping count with the lowest population rare risk, searching every stop from
zero through the cap. The curves connect these numerical minima; they are not
trajectories of a single trained model. Selection uses population risk averaged
over training randomness, rather than validation observations.

We repeat the search with every second positive grid point and check all 244
selected endpoints against a separate full tensor moment operator constructed
by enumerating the four binary observations. The reproduction script requires
the maximum risk change on grid refinement to be below $2\times10^{-6}$ and the
maximum tensor replay discrepancy to be below $10^{-10}$. All selected rates
are interior to the grid. These checks assess numerical accuracy within the
stated search; the theorem supplies the separate comparison over all
nonnegative constant rates. The figure illustrates an existing theoretical
result and is not a prospective evaluation of its predictions.
